% =============================================================================
%  Notas de clase — Sesión 14: Apuntadores, Referencias y Arreglos
%  Programación Avanzada — Universidad Panamericana
%  Compilar: pdflatex sesion14_apuntadores_referencias_arreglos.tex
% =============================================================================
\documentclass[12pt, oneside]{article}
\usepackage{progavanzada}

\begin{document}

\portada{Sesión 14}{Apuntadores, Referencias y Arreglos}{2026}

\tableofcontents
\newpage

% =============================================================================
\section{La memoria del computador}
% =============================================================================

Todo programa necesita guardar datos mientras se ejecuta. Esos datos viven en
la \textbf{memoria RAM}, que puede imaginarse como un enorme arreglo de celdas
numeradas. Cada celda tiene una dirección única (su número de casilla) y
almacena exactamente un byte.

\begin{center}
\begin{tabular}{ccc}
  \toprule
  \textbf{Dirección} & \textbf{Contenido} & \textbf{Variable} \\
  \midrule
  \texttt{0x61ff08} & \texttt{16} & \texttt{int x = 16;} \\
  \texttt{0x61ff0c} & \texttt{?}  & (próxima variable) \\
  \bottomrule
\end{tabular}
\end{center}

Cuando declaras \cpp{int x = 16;}, el compilador reserva 4 bytes contiguos en
RAM, les asigna el valor 16 y les da el nombre \cpp{x}. El operador \cpp{\&}
(dirección-de) nos permite ver en cuál celda vive esa variable.

\begin{codigo}[caption={Direccion de una variable}]
int x = 16;
cout << "Valor:     " << x  << endl;  // 16
cout << "Direccion: " << &x << endl;  // 0x61ff08  (varia cada ejecucion)
cout << "Tamano:    " << sizeof(x) << " bytes" << endl;  // 4
\end{codigo}

\begin{consejo}
  Las direcciones se muestran en \textbf{hexadecimal} (base 16, prefijo
  \texttt{0x}). El valor exacto cambia en cada ejecución porque el sistema
  operativo decide dónde colocar las variables (\textit{ASLR — Address Space
  Layout Randomization}).
\end{consejo}

\subsection{Tamaño de los tipos fundamentales}

\begin{center}
\begin{tabular}{lcc}
  \toprule
  \textbf{Tipo} & \textbf{Bytes (64-bit)} & \textbf{Rango típico} \\
  \midrule
  \texttt{char}   & 1 & $-128$ a $127$ \\
  \texttt{int}    & 4 & $\approx \pm 2\,000\,000\,000$ \\
  \texttt{long}   & 8 & $\approx \pm 9 \times 10^{18}$ \\
  \texttt{float}  & 4 & $\approx \pm 3.4 \times 10^{38}$ \\
  \texttt{double} & 8 & $\approx \pm 1.7 \times 10^{308}$ \\
  \texttt{bool}   & 1 & \texttt{true} / \texttt{false} \\
  \texttt{pointer}& 8 & toda la memoria virtual del proceso \\
  \bottomrule
\end{tabular}
\end{center}

\begin{tecnico}
  \textbf{Arquitectura Von Neumann (1945):} el principio de que los datos y el
  programa residen en la misma memoria es la base de todos los computadores
  modernos. Los apuntadores explotan directamente este modelo: son datos que
  apuntan a otros datos, o incluso al código mismo. \par\smallskip
  \textit{Referencia:} J. von Neumann, ``First Draft of a Report on the EDVAC'',
  1945. Disponible en: \url{https://archive.org/details/firstdraftofrepo00vonn}
\end{tecnico}

% =============================================================================
\section{Apuntadores}
% =============================================================================

Un \textbf{apuntador} es una variable cuyo valor es una \textit{dirección de
memoria}. En lugar de guardar un entero o un carácter, guarda \textit{dónde}
vive otro dato.

\begin{center}
\begin{tabular}{p{5cm} p{7cm}}
\toprule
\textbf{Variable normal} & \textbf{Apuntador} \\
\midrule
\begin{minipage}{5cm}
\begin{verbatim}
+---------+
|   16    | <- int x
+---------+
0x61ff08
\end{verbatim}
\end{minipage}
&
\begin{minipage}{7cm}
\begin{verbatim}
+--------------+
|  0x61ff08    | <- int* p
+------+-------+
       |
       v
+---------+
|   16    | <- int x
+---------+
\end{verbatim}
\end{minipage}
\\
\bottomrule
\end{tabular}
\end{center}

\subsection{Declaración}

\begin{codigo}
int* p;      // p es un apuntador a int
double* q;   // q es un apuntador a double
char* r;     // r es un apuntador a char
\end{codigo}

El tipo importa: el compilador necesita saber cuántos bytes leer al
desreferenciar. Un \cpp{int*} lee 4 bytes; un \cpp{double*} lee 8.

\subsection{Los dos operadores clave}

\begin{center}
\begin{tabular}{clp{7cm}}
  \toprule
  \textbf{Operador} & \textbf{Nombre} & \textbf{Qué hace} \\
  \midrule
  \cpp{\&x} & Dirección-de & Devuelve la dirección de memoria de \cpp{x} \\
  \cpp{*p}  & Desreferencia & Va a la dirección en \cpp{p} y lee/escribe el valor \\
  \bottomrule
\end{tabular}
\end{center}

\begin{advertencia}
  El símbolo \cpp{*} tiene \textbf{dos significados} según el contexto:\\
  \textbf{En una declaración} (\cpp{int* p;}): marca que \cpp{p} es un
  apuntador.\\
  \textbf{En una expresión} (\cpp{*p}): desreferencia — lee el valor apuntado.
  No los confundas.
\end{advertencia}

\begin{codigo}[caption={Apuntador: asignacion y desreferencia}]
int  x = 16;
int* p;          // declaro el apuntador (aun no apunta a nada valido)

p = &x;          // p guarda la direccion de x

cout << p   << endl;  // direccion: 0x61ff08
cout << *p  << endl;  // valor:     16
cout << &p  << endl;  // p tambien vive en algun lugar de memoria
\end{codigo}

\subsection{Modificar el original a través del apuntador}

\begin{codigo}[caption={Escritura a traves de apuntador}]
int  x = 16;
int* p = &x;

*p = 18;    // equivale exactamente a: x = 18;

cout << x;  // imprime 18
\end{codigo}

\begin{recuerda}
  \cpp{*p = 18} y \cpp{x = 18} son instrucciones \textbf{idénticas} cuando
  \cpp{p = \&x}. Ambas escriben el valor 18 en la misma celda de memoria.
\end{recuerda}

\subsection{Apuntador nulo}

Un apuntador que no apunta a nada válido debe inicializarse a \cpp{nullptr}.
Desreferenciar un apuntador nulo produce un \textit{segmentation fault} en
tiempo de ejecución.

\begin{codigo}
int* p = nullptr;   // apuntador nulo seguro

if (p != nullptr) {
    cout << *p;     // solo desreferenciamos si es seguro
}
\end{codigo}

\begin{tecnico}
  \textbf{Historia:} En C el valor nulo era la macro \texttt{NULL} (definida
  como \texttt{0} o \texttt{(void*)0}). C++11 introdujo \texttt{nullptr}, un
  literal con tipo propio (\texttt{std::nullptr\_t}), eliminando ambigüedades
  al sobrecargar funciones. Siempre usa \texttt{nullptr} en C++ moderno.\\
  \textit{Referencia:} Stroustrup, B. (2013). \textit{The C++ Programming
  Language}, 4.ª ed., §7.2.2.
\end{tecnico}

% =============================================================================
\section{Referencias}
% =============================================================================

Una \textbf{referencia} es un \textit{alias}: otro nombre para la misma
variable. No ocupa espacio adicional en memoria — comparte exactamente la
misma celda con la variable original.

\begin{codigo}[caption={Referencia: mismo espacio, otro nombre}]
int  x = 16;
int& y = x;    // y es alias de x

cout << &x << endl;  // 0x61ff08
cout << &y << endl;  // 0x61ff08  <- exactamente igual
\end{codigo}

\begin{consejo}
  La diferencia entre \cpp{\&} en una \textbf{declaración} y en una
  \textbf{expresión}:\\
  \cpp{int\& y = x;} → declaración de referencia (alias).\\
  \cpp{\&x} en una expresión → operador dirección-de (devuelve la dirección).
\end{consejo}

\subsection{Diferencias: Apuntador vs Referencia}

\begin{center}
\begin{tabular}{lll}
  \toprule
  \textbf{Característica} & \textbf{Apuntador (\cpp{int*})} & \textbf{Referencia (\cpp{int\&})} \\
  \midrule
  ¿Puede ser nulo?          & Sí (\cpp{nullptr})   & No — debe inicializarse \\
  ¿Puede reasignarse?       & Sí                   & No — fija tras la declaración \\
  ¿Ocupa memoria propia?    & Sí (8 bytes en 64b)  & No (alias) \\
  Sintaxis para leer valor  & \cpp{*p}             & \cpp{r} (directo) \\
  Uso típico                & Memoria dinámica     & Paso por referencia \\
  \bottomrule
\end{tabular}
\end{center}

\begin{tecnico}
  \textbf{¿Por qué existen las referencias?} Bjarne Stroustrup las introdujo
  en C++ (ausentes en C) precisamente para soportar el paso por referencia en
  funciones con una sintaxis más segura y clara que los apuntadores crudos.\\
  \textit{Referencia:} Stroustrup, B. (1994). \textit{The Design and Evolution
  of C++}, §3.7.
\end{tecnico}

% =============================================================================
\section{Arreglos}
% =============================================================================

Un arreglo es un bloque de celdas \textbf{contiguas} en memoria, todas del
mismo tipo. Para \cpp{int arr[5]}, el compilador reserva exactamente $5 \times
4 = 20$ bytes consecutivos:

\begin{center}
\begin{tabular}{ccccc}
  \toprule
  \texttt{arr[0]} & \texttt{arr[1]} & \texttt{arr[2]} & \texttt{arr[3]} & \texttt{arr[4]} \\
  \midrule
  10 & 20 & 30 & 40 & 50 \\
  \texttt{0x100} & \texttt{0x104} & \texttt{0x108} & \texttt{0x10C} & \texttt{0x110} \\
  \bottomrule
\end{tabular}
\end{center}

El nombre del arreglo es un \textbf{apuntador constante} al primer elemento:
\cpp{arr} $\equiv$ \cpp{\&arr[0]}.

\subsection{Aritmética de apuntadores}

Sumar un entero a un apuntador no suma bytes: suma \textbf{múltiplos del
tamaño del tipo}.

\begin{codigo}[caption={Aritmetica de apuntadores}]
int arr[5] = {10, 20, 30, 40, 50};

cout << arr     << endl;  // 0x100  (direccion del primer elemento)
cout << arr + 1 << endl;  // 0x104  (+4 bytes, no +1)
cout << arr + 2 << endl;  // 0x108  (+8 bytes)

cout << *arr       << endl;  // 10  (arr[0])
cout << *(arr + 1) << endl;  // 20  (arr[1])
cout << *(arr + 2) << endl;  // 30  (arr[2])
\end{codigo}

\subsection{Equivalencia fundamental}

Para cualquier arreglo, las dos notaciones son \textbf{idénticas} en el
código máquina generado:

\begin{center}
  \Large $\texttt{arr[i]} \;\equiv\; \texttt{*(arr + i)}$
\end{center}

El compilador transforma internamente \cpp{arr[i]} en \cpp{*(arr + i)}.
Los corchetes son azúcar sintáctica.

\begin{codigo}[caption={Dos formas equivalentes de recorrer un arreglo}]
int arr[5] = {10, 20, 30, 40, 50};

for (int i = 0; i < 5; i++) {
    cout << arr[i]       // notacion clasica
         << " == "
         << *(arr + i)   // notacion con apuntador
         << endl;
}
\end{codigo}

\begin{tecnico}
  \textbf{Dato curioso:} dado que \cpp{arr[i]} $\equiv$ \cpp{*(arr + i)} y
  la suma es conmutativa, \cpp{arr[i]} $\equiv$ \cpp{i[arr]} es una expresión
  válida en C y C++. Es un truco de entrevistas técnicas — jamás lo uses en
  código real.
\end{tecnico}

% =============================================================================
\section{Comportamiento indefinido: acceso fuera de límites}
% =============================================================================

C++ \textbf{no verifica} los límites de un arreglo en tiempo de ejecución.
Acceder fuera es \textit{undefined behavior} (UB): el resultado es impredecible.

\begin{codigo}
int arr[5] = {10, 20, 30, 40, 50};

cout << arr[5];   // UB: lee memoria fuera del arreglo
cout << arr[-1];  // UB: idem hacia atras
\end{codigo}

\begin{advertencia}
  El comportamiento indefinido no significa que el programa ``fallará''. Puede
  imprimir basura, devolver datos de otra variable, no fallar en la máquina
  del desarrollador y sí en la del usuario, o permitir que un atacante
  controle el flujo del programa. \textbf{Es la fuente de muchas
  vulnerabilidades de seguridad críticas.}
\end{advertencia}

\begin{tecnico}
  \textbf{El Gusano Morris (1988):} el primer malware de gran escala en
  Internet explotó un \textit{buffer overflow} (acceso fuera de límites en un
  arreglo) en el daemon \texttt{fingerd} de Unix. Infectó $\sim$6\,000
  máquinas, el 10\% del Internet de la época, y le costó a su autor la primera
  condena bajo la \textit{Computer Fraud and Abuse Act}.\\
  \textit{Referencia:} Spafford, E. (1989). ``The Internet Worm Program: An
  Analysis''. \textit{ACM SIGCOMM Computer Communication Review}, 19(1),
  17--57. \url{https://doi.org/10.1145/66093.66095}
\end{tecnico}

Para acceso seguro, usa \cpp{std::vector} con el método \cpp{.at(i)}, que
lanza una excepción si el índice está fuera de rango:

\begin{codigo}
#include <vector>
vector<int> v = {10, 20, 30, 40, 50};
cout << v.at(2);   // seguro: lanza std::out_of_range si i >= size
cout << v.at(9);   // lanza excepcion -- no UB
\end{codigo}

% =============================================================================
\section{Relación arreglo / apuntador}
% =============================================================================

\subsection{Array decay}

Cuando usas el nombre de un arreglo en casi cualquier expresión, se convierte
automáticamente en un apuntador al primer elemento. Esta conversión se llama
\textbf{\textit{array decay}} (decaimiento de arreglo).

\begin{codigo}
int arr[5] = {10, 20, 30, 40, 50};
int* p = arr;   // equivale a: int* p = &arr[0];

// p y arr se comportan igual para indexar...
cout << p[2];   // 30  <- igual que arr[2]

// ...pero sizeof los distingue:
cout << sizeof(arr) << endl;  // 20  (5 elementos x 4 bytes)
cout << sizeof(p)   << endl;  //  8  (tamano de un apuntador en 64-bit)
\end{codigo}

\begin{recuerda}
  \textbf{Consecuencia práctica del array decay:} cuando pasas un arreglo a
  una función, la función solo recibe un apuntador; pierde la información del
  tamaño. Por eso las funciones que operan sobre arreglos siempre reciben el
  tamaño como parámetro adicional.
\end{recuerda}

% =============================================================================
\section{Para saber más}
% =============================================================================

\subsection{Lecturas recomendadas}

\begin{itemize}
  \item \textbf{cppreference — Pointer declaration:}
        \url{https://en.cppreference.com/w/cpp/language/pointer}
  \item \textbf{cppreference — Reference declaration:}
        \url{https://en.cppreference.com/w/cpp/language/reference}
  \item \textbf{Kernighan \& Ritchie,} \textit{The C Programming Language},
        2.ª ed. (1988), Capítulo 5: Pointers and Arrays. El capítulo que
        definió cómo se enseñan los apuntadores durante décadas.
  \item \textbf{Stroustrup, B.} \textit{A Tour of C++}, 3.ª ed. (2022),
        \S1.7--1.8. Visión moderna y concisa de apuntadores en C++.
  \item \textbf{C++ Core Guidelines — Bounds safety:}
        \url{https://isocpp.github.io/CppCoreGuidelines/CppCoreGuidelines\#S-bounds}
\end{itemize}

\subsection{Contexto histórico}

\begin{tecnico}[Historia de los apuntadores en C]
  Los apuntadores llegaron a C de mano de \textbf{Dennis Ritchie} en los
  laboratorios Bell (1972), inspirados en el lenguaje B de Ken Thompson,
  que a su vez derivaba del BCPL de Martin Richards (1967). La motivación
  era práctica: escribir Unix de manera portable y eficiente sin ensamblador.
  \par\smallskip
  Ritchie recibió el Premio Turing en 1983 junto con Thompson ``por el
  desarrollo del la teoría de los sistemas operativos genéricos, y en
  particular por la implementación del sistema operativo UNIX''.
  \par\smallskip
  \textit{Referencia:} Ritchie, D. M. (1993). ``The Development of the C
  Language''. \textit{ACM SIGPLAN Notices}, 28(3), 201--208.
  \url{https://dl.acm.org/doi/10.1145/154766.155580}
\end{tecnico}

% =============================================================================
\section{Resumen}
% =============================================================================

\begin{recuerda}
\begin{center}
\begin{tabular}{lll}
  \toprule
  \textbf{Concepto} & \textbf{Símbolo} & \textbf{Qué hace} \\
  \midrule
  Dirección-de   & \cpp{\&x}         & Dirección de memoria de \cpp{x} \\
  Apuntador      & \cpp{int* p}      & Variable que guarda una dirección \\
  Desreferencia  & \cpp{*p}          & Lee/escribe el valor en la dirección \\
  Referencia     & \cpp{int\& r = x} & Alias: mismo espacio, otro nombre \\
  Array decay    & \cpp{arr} en expr.& Se convierte en \cpp{\&arr[0]} \\
  Aritmética     & \cpp{p + i}       & Avanza \cpp{i*sizeof(tipo)} bytes \\
  Equivalencia   & \cpp{arr[i]}      & Idéntico a \cpp{*(arr+i)} \\
  \bottomrule
\end{tabular}
\end{center}
\end{recuerda}

\end{document}
